\section{Discussion and Limitations}\label{sec:discussion}
The results in Section~\ref{sec:results} establish what the system does
reliably and where its evidence stops. This section interprets them, explains
the local actor-inference boundary, relates \textit{Boyun} to
\textit{MiroFish}, and sets out the work needed before any predictive claim
could be made.

\subsection{What the Results Show}
\begin{revised}
Taken together, the mechanism checks answer RQ1 and RQ2 within their stated
scope: arithmetic, timing, accounting, and saved-decision replay agree with
specified rules. Changing individual mechanisms changes results in controlled
ways. In particular, the offline preference test demonstrates that changing
numeric weights affects its policy, whereas rewriting persona prose alone does
not. That test does not establish how a learned actor model interprets either.

Extending those checks to a larger setting, local \textit{Laya} inference
carries all 24 token-launch paths through 336 hours. This verifies the ability
to record and settle repeated actor decisions without a per-call API fee.
It does not validate the regime hypotheses or the trading behavior. Realized
cash favors staged exit in the observed pairs, yet including inventory quotes
changes that comparison in six pairs. The metric selected therefore affects
what a reader may conclude.

Apart from trading, the four local lending worlds distinguish a numerical
constraint from behavior that happens to reach it. A borrower cap limits the
maximum possible loan, but reaching that cap depends on model choices. Likewise,
zero observed bad debt is conditional on the tested prices, liquidation rules,
liquidator resources, and depth. Reports should identify these assumptions
rather than treating either result as an unconditional property of a launch.
\end{revised}

\subsection{Self-Hosted Actor Inference}
\begin{revised}
The local execution path addresses a practical privacy requirement:
world observations can be scored on the user's hardware instead of being sent
to a hosted actor API. This narrows a data-transfer boundary; it is not a proof
of system-wide confidentiality. External data services, optional commitments,
artifact sharing, and host access remain separate considerations. Zero API
fee likewise excludes hardware, electricity, setup, and maintenance costs.

Within that local path, \textit{Laya}'s limitations are consequential.
Its context is limited to 1024 tokens, observations that would be truncated are
rejected, and its option scores are uncalibrated \cite{layamodel}. The recorded
trading distributions show limited regime separation, while lending choices
frequently select borrowing under both tested rate curves. These patterns
should motivate validation against actual behavior, not be presented as
measured investor or borrower demand.

Looking beyond the current experiments, other open-weight typed models, such
as \textit{Kev}, are candidates for comparison \cite{kevrepo}. They have not
been evaluated in these case studies, and compatibility alone would not imply
behavioral quality. Multiple seeds, measured demand, and matched evaluations
are stronger next steps than claiming predictive value from local inference
or from agreement between models.
\end{revised}

\subsection{Relation to \textit{MiroFish}}
The comparison with \textit{MiroFish} suggests a division of labor rather than
a contest. \textit{MiroFish} produces rich, persona-specific discourse but no
financial consequences; \textit{Boyun} produces settled, replayable financial
consequences but only pre-written messages. A natural next step is to let
language-model actors write \textit{Boyun}'s messages while the engine keeps
control of delivery, timing, and settlement, so that what an actor says
is generated freely but what it can see and spend is still enforced. Doing so
would raise the cost per decision and introduce new questions about how text
should influence trading, which would need their own evaluation.

\subsection{Toward Predictive Validation}
None of the tests in this paper measures forecasting accuracy, and
Figure~\ref{fig:validation} outlines what such a test would require. A study
must first fix one measurable outcome, its resolution source, and its
horizon, and freeze the world-construction procedure, models, and evidence
cutoff before any forecast is issued. Forecasts must then be dated and kept,
including failed runs, and scored against resolved outcomes using
Eqs.~\eqref{eq:brier} and~\eqref{eq:logloss}, alongside simple baselines such
as a historical base rate and a single-model forecast given the same
evidence. Cases that share a time period or narrative are not independent and
must be treated as such when estimating uncertainty. A separate user study
would be needed to test whether \textit{Boyun}'s reports improve decisions,
which is a different question from whether its probabilities are accurate.

\begin{figure}[t]
\centering
\begin{tikzpicture}[font=\scriptsize,
 vstep/.style={draw=teal!70,fill=wash,rounded corners=2pt,align=center,text width=3.4cm,minimum height=0.85cm,inner sep=3pt},
 arr/.style={-{Latex[length=1.6mm]},thick,draw=teal}]
\node[vstep] (a) at (0,0) {Fix outcome, source,\\and horizon};
\node[vstep] (b) at (4.1,0) {Freeze construction,\\models, and evidence cutoff};
\node[vstep] (c) at (4.1,-1.35) {Issue and keep dated\\forecasts, incl.\ failures};
\node[vstep] (d) at (0,-1.35) {Resolve outcomes from\\the declared source};
\node[vstep] (e) at (0,-2.7) {Score against baselines;\\report coverage and cost};
\node[vstep] (f) at (4.1,-2.7) {Limit claims to\\the scored evidence};
\draw[arr] (a)--(b); \draw[arr] (b)--(c); \draw[arr] (c)--(d);
\draw[arr] (d)--(e); \draw[arr] (e)--(f);
\end{tikzpicture}
\caption{Proposed prospective validation protocol. This is future work,
not an experiment reported in this paper.}
\label{fig:validation}
\end{figure}

\subsection{Limitations}
\begin{revised}
Several limitations apply to the full evaluation. The case results use one
seed per plan, synthetic actor balances and preferences, and declared event or
price schedules. They do not represent independent observations of human
behavior. In \textit{MACRON}, near-uniform option scores weaken regime
separation, and upfront spending alone triggers the loss-budget alert in the
higher-upfront plans. Founder trades are partly policy-driven, so model choice
cannot explain every cash difference.

On the market side, a simplified constant-product pool omits order books,
transaction ordering, front-running, outages, and venue-specific execution
details. On the lending side, assumed liquidator cash and exit depth determine
which liquidations can settle; historical hourly closes do not test intrahour
price shocks. The competing venue's rates are fixed at one snapshot and campaign
reach has no measured conversion. These conditions bound the scope of every
financial conclusion.

Finally, messages attached to actions are pre-written, giving \textit{Boyun}
a narrower social layer than open-ended discourse. The \textit{MiroFish}
comparison matches the brief, but actor models, activation rules, and recovery
procedures differ. It compares recorded capabilities rather than forecasting
accuracy or general superiority.
\end{revised}

\subsection{Responsible Use}
These limits also shape how \textit{Boyun} should and should not be used.
Its outputs are conditional explorations of stated assumptions, not
financial advice or price forecasts. Because the case study concerns paid
promotion of a speculative token, the same tools could in principle be used to
rehearse campaigns that mislead buyers; reports therefore state their
assumptions and limits explicitly, and they should not be presented to
investors as evidence of likely returns. World files and reports can contain
private strategy, so authors should share them deliberately. Local actor
inference avoids sending observations to an actor API; external services and
publication remain separate disclosure decisions.
